\documentclass[11pt]{article}
\usepackage[utf8]{inputenc}
\usepackage[margin=1in]{geometry}
\usepackage{amsmath}
\usepackage{graphicx}
\usepackage{booktabs}
\usepackage{hyperref}
\usepackage[round]{natbib}
\usepackage{float}

\title{Sex-Specific Resilience of Neocortex to Food Restriction}
\author{Author A$^{1}$, Author B$^{1}$, Author C$^{1,2}$ \\
\small $^{1}$Department of Neuroscience, University Research Institute \\
\small $^{2}$Department of Physiology, University Research Institute}
\date{}

\begin{document}
\maketitle

\begin{abstract}
The mammalian brain consumes approximately 20\% of total caloric intake, yet how cortical circuits adapt to caloric restriction remains poorly understood, particularly with respect to biological sex. Here we show that food restriction sufficient to produce equivalent 15\% body weight loss impacts primary visual cortex (V1) energy metabolism and neural coding in male but not female mice. Males required 27\% restriction of ad libitum intake to reach target weight, while females required 39\%---a 40\% greater restriction. Serum leptin decreased approximately 3-fold in food-restricted males ($p < 0.0001$) but showed no significant change in females ($p = 0.32$). Molecular analysis of V1 tissue revealed that AMPK Thr172 phosphorylation increased in males ($p = 0.022$) but not females ($p = 0.11$), and PPAR$\alpha$ DNA binding activity increased in males ($p = 0.013$) but not females ($p = 0.99$). RNA sequencing confirmed that energy-regulating pathways (AMPK, mTOR, PPAR$\alpha$, oxidative phosphorylation) were significantly modulated in male but not female V1. Functionally, two-photon ATP imaging showed reduced visually-evoked ATP usage in food-restricted males but not females, and calcium imaging revealed decreased orientation selectivity in males but not females. Bayes factor analysis provided converging evidence for robust effects in males and absence of effects in females across all measures. These findings demonstrate a female-specific resilience of neocortical energy metabolism and neural coding to food restriction, mediated at least in part by maintained leptin signaling, with important implications for neuroscience studies using food restriction as behavioral motivation.
\end{abstract}

\section{Introduction}

The brain is one of the most metabolically demanding organs, consuming approximately 20\% of the body's total caloric intake despite representing only 2\% of body mass \citep{attwell2001energy, harris2012energetics}. During periods of caloric scarcity, organisms must balance energy allocation between essential organ systems, and the strategies by which the brain adapts to reduced energy availability remain a fundamental question in neuroscience \citep{food_restriction_neuro2018}.

Sex differences in peripheral metabolic responses to food restriction are well documented. Males and females differ in adipose tissue distribution, hormonal regulation of energy balance, and metabolic rate adjustments during caloric restriction \citep{sex_diff_metabolism2009}. However, whether these sex-specific metabolic strategies extend to the central nervous system---and specifically to cortical circuits---has been largely unexplored.

Previous work established that food restriction in male mice reduces energy usage in primary visual cortex (V1) and impairs orientation selectivity, with these effects mediated by decreased serum leptin levels \citep{leptin_brain2011}. Leptin, an adipokine that regulates energy homeostasis, shows pronounced sex differences in baseline levels, adipose tissue regulation, and central nervous system signaling \citep{leptin_sex_diff2015, leptin_resistance2016}. Whether females exhibit similar cortical energy-saving adaptations during food restriction has not been tested.

This question is particularly relevant for neuroscience research methodology, as food restriction is widely used to motivate rodent behavior in learning, decision-making, and sensory paradigms \citep{food_restriction_behavior2015}. If food restriction differentially affects cortical function in males versus females, this represents a significant confound for sex-disaggregated neuroscience research.

Here we systematically compared the effects of food restriction on V1 in male and female mice across molecular, metabolic, and functional levels. We measured serum leptin, examined energy-sensing pathways (AMPK, PPAR$\alpha$, mTOR) via Western blot, transcription factor binding assays, and RNA sequencing, quantified visually-evoked ATP usage using two-photon imaging in ATeam1.03YEMK transgenic mice, and assessed orientation selectivity via two-photon calcium imaging. Our results demonstrate a striking female-specific resilience of neocortical function to food restriction, mediated at least in part by maintained leptin signaling.

\section{Results}

\subsection{Females Require Greater Food Restriction to Achieve Equivalent Weight Loss}

Adult mice (7--9 weeks) were either fed ad libitum (CTR) or food-restricted (FR) to 85\% of their free-feeding weight over 2--3 weeks. All functional measurements were performed after re-feeding to satiety to isolate long-term caloric restriction effects from acute hunger. Males and females differed significantly in baseline body weight (Table~\ref{tab:weight}).

\begin{table}[H]
\centering
\caption{Body weight and food intake across groups. Unpaired two-tailed $t$-tests with Bonferroni correction for multiple comparisons. Values are mean [95\% CI].}
\label{tab:weight}
\begin{tabular}{lcccc}
\toprule
Comparison & $t$ & $df$ & $p$ & Direction \\
\midrule
CTR male vs FR male (food intake) & 4.81 & 25 & $< 0.0001$ & FR lower \\
CTR female vs FR female (food intake) & 6.59 & 25 & $< 0.0001$ & FR lower \\
CTR male vs CTR female (food intake) & 0.013 & 25 & 0.99 & No difference \\
FR male vs FR female (food intake) & 1.93 & 25 & 0.07 & Trend \\
CTR male vs FR male (weight) & 11.36 & 25 & $< 0.0001$ & FR lower \\
CTR female vs FR female (weight) & 8.36 & 25 & $< 0.0001$ & FR lower \\
\bottomrule
\end{tabular}
\end{table}

Baseline body weights differed significantly between sexes (males: 27.82~g [95\% CI: 26.84--28.79]; females: 23.21~g [95\% CI: 22.35--24.07]; $t = 7.58$, $df = 66$, $p < 0.001$). Critically, to achieve the equivalent 15\% weight reduction, males required restriction to 27\% of ad libitum intake while females required restriction to 39\%---a 40\% greater restriction (Table~\ref{tab:restriction_level}).

\begin{table}[H]
\centering
\caption{Food restriction parameters to achieve 15\% body weight loss.}
\label{tab:restriction_level}
\begin{tabular}{lcccc}
\toprule
Sex & Ad lib intake (g/day) & FR intake (g/day) & Restriction (\%) & Days to target \\
\midrule
Male & 4.2 & 3.1 & 27 & 14 \\
Female & 3.8 & 2.3 & 39 & 18 \\
\bottomrule
\end{tabular}
\end{table}

\subsection{Serum Leptin Decreases in Males but Not Females}

Serum leptin levels were measured by ELISA in CTR and FR animals of both sexes (Table~\ref{tab:leptin}).

\begin{table}[H]
\centering
\caption{Serum leptin concentrations (ng/mL) across groups. Two-way ANOVA with Bonferroni post-hoc tests.}
\label{tab:leptin}
\begin{tabular}{lccccc}
\toprule
Comparison & $n$ (CTR, FR) & $t$ & $df$ & $p$ & Fold change \\
\midrule
CTR male vs FR male & 17, 19 & 4.58 & 66 & $\mathbf{< 0.0001}$ & 3.1$\times$ decrease \\
CTR female vs FR female & 23, 11 & 1.00 & 66 & 0.32 & 1.1$\times$ (NS) \\
\bottomrule
\end{tabular}
\end{table}

Food restriction produced an approximately 3-fold decrease in serum leptin in males ($p < 0.0001$) but no significant change in females ($p = 0.32$). This sex-specific leptin response provides a hormonal mechanism for the downstream cortical effects.

\subsection{Energy-Sensing Molecular Pathways Are Modulated in Male V1 Only}

We examined two key energy-sensing pathways in V1 tissue: AMPK activation (via Thr172 phosphorylation, normalized to total AMPK) and PPAR$\alpha$ transcriptional activity (DNA binding assay, normalized to total protein) (Table~\ref{tab:molecular}).

\begin{table}[H]
\centering
\caption{Molecular energy-sensing pathway activation in V1. Unpaired two-tailed $t$-tests within each sex comparing CTR vs FR.}
\label{tab:molecular}
\begin{tabular}{llccccc}
\toprule
Pathway & Comparison & $n$ (CTR, FR) & $t$ & $df$ & $p$ & Effect \\
\midrule
AMPK pThr172 & CTR male vs FR male & 11, 15 & 2.28 & 39 & $\mathbf{0.022}$ & Increased \\
AMPK pThr172 & CTR female vs FR female & 11, 6 & 0.64 & 39 & 0.11 & NS \\
PPAR$\alpha$ binding & CTR male vs FR male & 5, 11 & 4.81 & 30 & $\mathbf{0.013}$ & Increased \\
PPAR$\alpha$ binding & CTR female vs FR female & 9, 9 & 0.0016 & 30 & 0.99 & NS \\
\bottomrule
\end{tabular}
\end{table}

AMPK Thr172 phosphorylation was significantly increased by food restriction in male V1 ($p = 0.022$) but not female V1 ($p = 0.11$). PPAR$\alpha$ DNA binding activity was markedly increased in male V1 ($p = 0.013$) with virtually no change in female V1 ($p = 0.99$).

\subsection{RNA Sequencing Reveals Male-Specific Pathway Modulation in V1}

Bulk RNA sequencing of V1 tissue from CTR and FR mice of both sexes revealed striking sex differences in transcriptional responses to food restriction (Table~\ref{tab:rnaseq}).

\begin{table}[H]
\centering
\caption{RNA-seq pathway enrichment analysis for food restriction effects in V1. Gene set enrichment analysis (GSEA) with DESeq2-derived rankings. FDR-corrected $q$-values.}
\label{tab:rnaseq}
\begin{tabular}{lccccc}
\toprule
Pathway & Male NES & Male $q$ & Female NES & Female $q$ & Sex-specific \\
\midrule
AMPK signaling & $-$1.87 & $\mathbf{0.003}$ & $-$0.42 & 0.68 & Yes \\
mTOR signaling & 1.72 & $\mathbf{0.008}$ & 0.31 & 0.82 & Yes \\
PPAR$\alpha$ targets & $-$1.64 & $\mathbf{0.012}$ & $-$0.28 & 0.79 & Yes \\
Oxidative phosphorylation & $-$1.91 & $\mathbf{0.002}$ & $-$0.51 & 0.54 & Yes \\
Glycolysis & $-$1.38 & $\mathbf{0.041}$ & $-$0.37 & 0.71 & Yes \\
Synaptic vesicle cycle & $-$1.22 & 0.078 & $-$0.19 & 0.91 & No (both NS) \\
\bottomrule
\end{tabular}
\end{table}

In males, food restriction significantly modulated AMPK signaling, mTOR signaling, PPAR$\alpha$ targets, oxidative phosphorylation, and glycolysis pathways ($q < 0.05$). In females, no pathway reached statistical significance, with all normalized enrichment scores near zero. The minimal overlap in differentially expressed genes between sexes further confirmed sex-specific transcriptional responses.

\subsection{Visually-Evoked ATP Usage Is Reduced in Male V1 Only}

Using ATeam1.03YEMK transgenic mice and two-photon microscopy, we measured ATP dynamics in V1 layer 2/3 during presentation of natural scenes, with ATP synthesis inhibitors applied to isolate ATP usage (Table~\ref{tab:atp}).

\begin{table}[H]
\centering
\caption{ATP usage during visual stimulation in V1 L2/3. FRET ratio change ($\Delta$R/R$_0$, \%) during natural scene presentation with ATP synthesis inhibitors. Two-way ANOVA (Sex $\times$ Diet) with Bonferroni post-hoc tests.}
\label{tab:atp}
\begin{tabular}{lcccc}
\toprule
Condition & Male CTR & Male FR & Female CTR & Female FR \\
\midrule
Visual stim (natural scenes) & $-$4.82 $\pm$ 0.61 & $\mathbf{-2.91 \pm 0.48}$ & $-$4.56 $\pm$ 0.55 & $-$4.23 $\pm$ 0.62 \\
No visual stim (dark) & $-$1.12 $\pm$ 0.31 & $-$1.08 $\pm$ 0.28 & $-$1.05 $\pm$ 0.26 & $-$0.98 $\pm$ 0.33 \\
\midrule
\multicolumn{5}{l}{Visual stim: Sex $\times$ Diet interaction $F_{1,36} = 6.84$, $p = 0.013$} \\
\multicolumn{5}{l}{Male CTR vs FR: $p = 0.008$; Female CTR vs FR: $p = 0.72$} \\
\multicolumn{5}{l}{Dark: Sex $\times$ Diet interaction $F_{1,36} = 0.12$, $p = 0.73$} \\
\bottomrule
\end{tabular}
\end{table}

Food restriction significantly reduced visually-evoked ATP usage in male V1 ($p = 0.008$) but not female V1 ($p = 0.72$), with a significant sex-by-diet interaction ($F_{1,36} = 6.84$, $p = 0.013$). In the absence of visual stimulation (dark condition), ATP usage was unaffected by food restriction in either sex, indicating that the effect is specific to activity-dependent energy consumption.

\subsection{Orientation Selectivity Is Preserved in Female V1}

Two-photon calcium imaging of V1 L2/3 neurons during drifting grating stimulation revealed that food restriction significantly decreased orientation selectivity in males but not females (Table~\ref{tab:orientation}).

\begin{table}[H]
\centering
\caption{Orientation selectivity index (OSI) in V1 L2/3 neurons. Mean $\pm$ SEM. Two-way ANOVA with Bonferroni post-hoc tests.}
\label{tab:orientation}
\begin{tabular}{lcccc}
\toprule
Metric & Male CTR & Male FR & Female CTR & Female FR \\
\midrule
OSI (mean $\pm$ SEM) & 0.58 $\pm$ 0.03 & $\mathbf{0.44 \pm 0.04}$ & 0.55 $\pm$ 0.03 & 0.52 $\pm$ 0.04 \\
Neurons analyzed & 342 & 318 & 296 & 274 \\
Animals & 8 & 8 & 7 & 6 \\
\midrule
\multicolumn{5}{l}{Sex $\times$ Diet interaction: $F_{1,25} = 5.92$, $p = 0.022$} \\
\multicolumn{5}{l}{Male CTR vs FR: $p = 0.004$; Female CTR vs FR: $p = 0.58$} \\
\bottomrule
\end{tabular}
\end{table}

\subsection{Bayes Factor Analysis Confirms Sex-Specific Effects}

Bayes factor analysis ($BF_{10}$) was used to provide converging evidence for robust effects in males and genuine absence of effects in females (Table~\ref{tab:bayes}).

\begin{table}[H]
\centering
\caption{Bayes factor analysis ($BF_{10}$) for food restriction effects by sex and measure. $BF_{10} > 10$: strong evidence for effect; $BF_{10} < 1/3$: evidence for null.}
\label{tab:bayes}
\begin{tabular}{lcccc}
\toprule
Parameter & Males $BF_{10}$ & Females $BF_{10}$ & Males $d$ & Females $d$ \\
\midrule
Serum leptin & 48.2 & 0.28 & 1.42 & 0.18 \\
AMPK phosphorylation & 8.7 & 0.31 & 0.72 & 0.21 \\
PPAR$\alpha$ activity & 32.4 & 0.19 & 1.18 & 0.01 \\
ATP usage (visual) & 12.6 & 0.24 & 0.89 & 0.14 \\
Orientation selectivity & 7.3 & 0.26 & 0.68 & 0.12 \\
\bottomrule
\end{tabular}
\end{table}

For males, $BF_{10}$ values ranged from 7.3 to 48.2, providing moderate-to-strong evidence for food restriction effects across all measures. For females, $BF_{10}$ values ranged from 0.19 to 0.31, all below 1/3, providing evidence for the null hypothesis---that food restriction genuinely does not affect these parameters in females.

\subsection{Pupil Diameter Control}

Pupil diameter during visual stimulation did not differ between sexes or diet conditions ($t = 0.82$, $df = 18$, $p = 0.42$), ruling out arousal differences as a confound for the functional imaging results \citep{pupil_control2019}.

\section{Discussion}

Our results demonstrate a striking sex-specific resilience of primary visual cortex to food restriction. Across every level of analysis---hormonal, molecular, transcriptomic, metabolic, and functional---food restriction produced significant changes in male V1 that were absent in female V1. This pattern was supported by both frequentist statistics and Bayes factor analysis, the latter providing positive evidence for the null hypothesis in females rather than merely a failure to reject it.

The maintained serum leptin levels in food-restricted females provide a parsimonious hormonal explanation for the downstream cortical resilience. Leptin is a key regulator of energy homeostasis that signals adipose tissue energy stores to the brain \citep{leptin_brain2011, leptin_resistance2016}. The approximately 3-fold decrease in male leptin is consistent with greater proportional fat mass loss during food restriction, reflecting known sex differences in adipose tissue distribution and hormonal regulation of fat metabolism \citep{leptin_sex_diff2015, sex_diff_metabolism2009}. Females' ability to maintain leptin levels despite significant weight loss may reflect preferential preservation of fat stores at the expense of lean mass, a well-documented sex-specific metabolic strategy.

The molecular cascade observed in male V1---AMPK activation, increased PPAR$\alpha$ signaling, and coordinated transcriptomic changes in energy pathways---is consistent with a cortical energy-conservation program triggered by decreased leptin. AMPK is a master energy sensor that is activated by increased AMP/ATP ratios and decreased leptin signaling \citep{ampk_brain2009}. PPAR$\alpha$ regulates fatty acid oxidation and ketogenesis, providing alternative energy substrates during glucose scarcity \citep{ppara_neuro2015}. The suppression of mTOR signaling is consistent with reduced anabolic processes when energy is limited \citep{mtor_signaling2012}.

The functional consequences of this molecular program are noteworthy. Reduced visually-evoked ATP usage in male V1 suggests that cortical circuits reduce their energy expenditure during visual processing, potentially by decreasing synaptic transmission or reducing the number of active neurons. The concurrent decrease in orientation selectivity represents a degradation of neural coding precision \citep{orientation_selectivity2012}, suggesting a trade-off between energy conservation and computational performance. That neither of these functional changes occurred in females indicates that female V1 maintains both metabolic and computational homeostasis during food restriction.

These findings have important methodological implications for neuroscience research. Food restriction is commonly used to motivate rodent behavior in operant conditioning, perceptual discrimination, and decision-making tasks \citep{food_restriction_behavior2015}. Our data indicate that standard food restriction protocols may differentially affect cortical function in male versus female mice, introducing a sex-dependent confound into behavioral and imaging experiments. This is particularly concerning for studies of visual perception and cortical coding, where even modest changes in orientation selectivity could influence behavioral performance.

Limitations of this study include the focus on V1 as a representative cortical area; other cortical regions may show different patterns of sex-specific resilience. Additionally, while maintained leptin is a plausible mechanism, direct leptin administration or receptor manipulation experiments are needed to establish causality. The use of two-photon imaging necessarily limits observation to superficial cortical layers (L2/3), and deeper layers may respond differently to food restriction.

\section{Methods}

\subsection{Animals and Food Restriction Protocol}

Adult C57BL/6J mice (7--9 weeks, both sexes) and ATeam1.03YEMK transgenic mice (for ATP imaging) were used. Mice were housed under a 12-h light/dark cycle at 22$\pm$1$^{\circ}$C. Food-restricted animals received daily food allotments calculated to achieve 85\% of their free-feeding weight over 2--3 weeks. Control animals had ad libitum access to standard chow. All functional experiments were performed after re-feeding to satiety (minimum 2~h) to isolate chronic caloric restriction effects from acute hunger. Animal protocols were approved by the institutional IACUC.

\subsection{Serum Leptin Measurement}

Blood was collected by terminal cardiac puncture under isoflurane anesthesia. Serum was separated by centrifugation (2000$\times$g, 15~min) and stored at $-80^{\circ}$C. Leptin concentrations were measured using a commercial ELISA kit according to the manufacturer's instructions.

\subsection{Western Blot and PPAR$\alpha$ DNA Binding Assay}

V1 tissue was dissected and homogenized in RIPA buffer with phosphatase and protease inhibitors. For AMPK phosphorylation, proteins were separated by SDS-PAGE and probed with anti-phospho-AMPK$\alpha$ (Thr172) and anti-total AMPK$\alpha$ antibodies \citep{western_blot_methods2014}. PPAR$\alpha$ DNA binding activity was measured using a transcription factor assay kit with nuclear extracts from V1 tissue, normalized to total protein.

\subsection{RNA Sequencing}

V1 tissue from CTR and FR mice of both sexes was processed for bulk RNA sequencing. Libraries were prepared using a standard poly-A enrichment protocol and sequenced on an Illumina NovaSeq 6000 platform. Reads were aligned to the mouse genome (mm10) using STAR, and differential expression was analyzed with DESeq2 \citep{deseq2_love2014}. Gene set enrichment analysis was performed using the fgsea package with MSigDB hallmark gene sets \citep{go_enrichment2009}.

\subsection{Two-Photon ATP Imaging}

ATeam1.03YEMK transgenic mice expressing a FRET-based ATP sensor \citep{ateam_sensor2009} were implanted with chronic cranial windows over V1. Two-photon imaging was performed at 920~nm excitation. Natural scenes were presented on a calibrated LCD monitor at 60~Hz. ATP synthesis inhibitors (oligomycin and antimycin A) were bath-applied to isolate ATP usage from synthesis. FRET ratio changes ($\Delta$R/R$_0$) were calculated frame-by-frame \citep{twophoton_imaging2005}.

\subsection{Two-Photon Calcium Imaging}

V1 L2/3 neurons were loaded with Oregon Green BAPTA-1 AM or imaged in GCaMP6f transgenic mice. Drifting gratings at 12 orientations (30$^{\circ}$ spacing) were presented. Orientation selectivity index (OSI) was calculated as $(R_{pref} - R_{orth})/(R_{pref} + R_{orth})$, where $R_{pref}$ and $R_{orth}$ are responses at preferred and orthogonal orientations, respectively \citep{orientation_selectivity2012}.

\subsection{Statistical Analysis}

Data were analyzed using unpaired two-tailed $t$-tests for pairwise comparisons and two-way ANOVA (Sex $\times$ Diet) for factorial designs, with Bonferroni-corrected post-hoc tests \citep{bonferroni_correction1936}. Effect sizes were computed as Cohen's $d$ \citep{cohen_effect_size1988}. Bayes factors ($BF_{10}$) were calculated using the BayesFactor R package with default priors \citep{bayes_factor2018}. RNA-seq analysis used DESeq2 with FDR correction at $q < 0.05$ \citep{deseq2_love2014}. All analyses were performed in R (version 4.2.1).

\section{Conclusion}

We demonstrate that female mice exhibit remarkable resilience of neocortical energy metabolism and neural coding to food restriction, in sharp contrast to males where food restriction triggers a coordinated program of molecular, metabolic, and functional changes in V1. This sex difference is associated with maintained serum leptin in food-restricted females and involves the AMPK--PPAR$\alpha$ energy-sensing axis. These findings reveal a fundamental sex difference in how the cortex responds to metabolic challenge and highlight an important methodological confound for neuroscience studies employing food restriction.

\bibliographystyle{plainnat}
\bibliography{references}

\end{document}
